\section{Representable functors}
\label{section:0.8}

\setcounter{subsection}{0}
\subsection{Representable functors}
\label{subsection:0.8.1}

\begin{env}[8.1.1]
\label{0.8.1.1}
We denote by $\mathrm{Ens}$ the category of sets. Let $C$ be a
category; for two objects $X$, $Y$ of $C$, we set
$h_X(Y)=\operatorname{Hom}(Y,X)$; for every morphism $u:Y\to Y'$ in $C$,
we denote by $h_X(u)$ the map $v\mapsto vu$ from
$\operatorname{Hom}(Y',X)$ to $\operatorname{Hom}(Y,X)$. It is immediate
that with these definitions, $h_X:C\to\mathrm{Ens}$ is a \emph{contravariant
functor}, that is, an object of the category, denoted
$\operatorname{Hom}(C^0,\mathrm{Ens})$, of covariant functors from the
category $C^0$, dual to the category $C$, to the category
$\mathrm{Ens}$ (T, 1.7, d) and [29]).
\end{env}

\begin{env}[8.1.2]
\label{0.8.1.2}
Now let $w:X\to X'$ be a morphism in $C$; for every $Y\in C$ and
every $v\in\operatorname{Hom}(Y,X)=h_X(Y)$, we have
$wv\in\operatorname{Hom}(Y,X')=h_{X'}(Y)$; denote by $h_w(Y)$
the map $v\mapsto wv$ from $h_X(Y)$ to $h_{X'}(Y)$. It is immediate
that for every morphism $u:Y\to Y'$ in $C$, the diagram
\[
  \xymatrix{
    h_X(Y')\ar[r]^{h_X(u)}\ar[d]_{h_w(Y')} &
    h_X(Y)\ar[d]^{h_w(Y)}\\
    h_{X'}(Y')\ar[r]_{h_{X'}(u)} &
    h_{X'}(Y)
  }
\]
is commutative; in other words, $h_w$ is a \emph{functorial morphism}
$h_X\to h_{X'}$ (T, 1.2), or again a morphism in the category
$\operatorname{Hom}(C^0,\mathrm{Ens})$ (T, 1.7, d)). The definitions of
$h_X$ and of $h_w$ therefore constitute the definition of a \emph{canonical
covariant functor}
\[
  h:C\to\operatorname{Hom}(C^0,\mathrm{Ens}),\qquad X\mapsto h_X.
  \tag{8.1.2.1}\label{0.8.1.2.1}
\]
\end{env}

\begin{env}[8.1.3]
\label{0.8.1.3}
Let $X$ be an object of $C$ and $F$ a contravariant functor from $C$ to
$\mathrm{Ens}$ (an object of $\operatorname{Hom}(C^0,\mathrm{Ens})$). Let
$g:h_X\to F$ be a \emph{functorial morphism}: for every $Y\in C$,
\oldpage[0\textsubscript{III}]{6}
$g(Y)$ is then a map $h_X(Y)\to F(Y)$ such that for every morphism
$u:Y\to Y'$ in $C$, the diagram
\[
  \xymatrix{
    h_X(Y')\ar[r]^{h_X(u)}\ar[d]_{g(Y')} &
    h_X(Y)\ar[d]^{g(Y)}\\
    F(Y')\ar[r]_{F(u)} & F(Y)
  }
  \tag{8.1.3.1}\label{0.8.1.3.1}
\]
is commutative. In particular, we have a map
$g(X):h_X(X)=\operatorname{Hom}(X,X)\to F(X)$, whence an element
\[
  \alpha(g)=(g(X))(1_X)\in F(X)
  \tag{8.1.3.2}\label{0.8.1.3.2}
\]
and consequently a canonical map
\[
  \alpha:\operatorname{Hom}(h_X,F)\to F(X).
  \tag{8.1.3.3}\label{0.8.1.3.3}
\]

Conversely, consider an element $\xi\in F(X)$; for every morphism
$v:Y\to X$ in $C$, $F(v)$ is a map $F(X)\to F(Y)$;
consider the map
\[
  v\mapsto(F(v))(\xi)
  \tag{8.1.3.4}\label{0.8.1.3.4}
\]
from $h_X(Y)$ to $F(Y)$; if we denote this map by
$(\beta(\xi))(Y)$, then
\[
  \beta(\xi):h_X\to F
  \tag{8.1.3.5}\label{0.8.1.3.5}
\]
is a \emph{functorial morphism}, since for every morphism
$u:Y\to Y'$ in $C$ we have
$(F(vu))(\xi)=(F(u)\circ F(v))(\xi)$, which verifies the commutativity of
(\hyperref[0.8.1.3.1]{8.1.3.1}) for $g=\beta(\xi)$. We have thus defined
a canonical map
\[
  \beta:F(X)\to\operatorname{Hom}(h_X,F).
  \tag{8.1.3.6}\label{0.8.1.3.6}
\]
\end{env}

\begin{proposition}[8.1.4]
\label{0.8.1.4}
The maps $\alpha$ and $\beta$ are bijections inverse to each
other.
\end{proposition}

Let us compute $\alpha(\beta(\xi))$ for $\xi\in F(X)$; for every $Y\in C$,
$(\beta(\xi))(Y)$ is the map
$g_1(Y):v\mapsto(F(v))(\xi)$ from $h_X(Y)$ to $F(Y)$. We therefore have
\[
  \alpha(\beta(\xi))=(g_1(X))(1_X)=(F(1_X))(\xi)
  =1_{F(X)}(\xi)=\xi.
\]
Let us now compute $\beta(\alpha(g))$ for
$g\in\operatorname{Hom}(h_X,F)$; for every $Y\in C$,
$(\beta(\alpha(g)))(Y)$ is the map
$v\mapsto(F(v))((g(X))(1_X))$; by virtue of the commutativity of
(\hyperref[0.8.1.3.1]{8.1.3.1}), this map is none other than
$v\mapsto(g(Y))((h_X(v))(1_X))=(g(Y))(v)$ by definition of $h_X(v)$;
in other words, it is equal to $g(Y)$, which proves the proposition.

\begin{env}[8.1.5]
\label{0.8.1.5}
Recall that a \emph{subcategory} $C'$ of a category $C$ is defined
by the condition that its objects are objects of $C$, and that if $X'$,
$Y'$ are two objects of $C'$, the set
$\operatorname{Hom}_{C'}(X',Y')$ of morphisms $X'\to Y'$ \emph{in
$C'$} is a subset of the set $\operatorname{Hom}_C(X',Y')$ of
morphisms $X'\to Y'$ \emph{in $C$}, the canonical map of
``composition of morphisms''
\[
  \operatorname{Hom}_{C'}(X',Y')\times\operatorname{Hom}_{C'}(Y',Z')
  \to\operatorname{Hom}_{C'}(X',Z')
\]
\oldpage[0\textsubscript{III}]{7}
being the restriction of the canonical map
\[
  \operatorname{Hom}_C(X',Y')\times\operatorname{Hom}_C(Y',Z')
  \to\operatorname{Hom}_C(X',Z').
\]

We say that $C'$ is a \emph{full subcategory} of $C$ if
$\operatorname{Hom}_{C'}(X',Y')=\operatorname{Hom}_C(X',Y')$ for every
pair of objects of $C'$. The subcategory $C''$ of $C$ consisting of the
objects of $C$ isomorphic to objects of $C'$ is then again a full
subcategory of $C$, \emph{equivalent} (T, 1.2) to $C'$, as is easily
verified.

A covariant functor $F:C_1\to C_2$ is said to be \emph{fully faithful} if,
for every pair of objects $X_1$, $Y_1$ of $C_1$, the map
$u\mapsto F(u)$ from $\operatorname{Hom}(X_1,Y_1)$ to
$\operatorname{Hom}(F(X_1),F(Y_1))$ is \emph{bijective}; this implies that
the subcategory $F(C_1)$ of $C_2$ is \emph{full}. Moreover, if two
objects $X_1$, $X'_1$ have the same image $X_2$, there exists a unique
isomorphism $u:X_1\to X'_1$ such that $F(u)=1_{X_2}$. For every object $X_2$
of $F(C_1)$, let then $G(X_2)$ be one of the objects $X_1$ of $C_1$ such that
$F(X_1)=X_2$ ($G$ being defined by means of the axiom of choice); for every
morphism $v:X_2\to Y_2$ in $F(C_1)$, $G(v)$ will be the unique morphism
$u:G(X_2)\to G(Y_2)$ such that $F(u)=v$; $G$ is then a \emph{functor} from
$F(C_1)$ to $C_1$; $FG$ is the identity functor on $F(C_1)$, and the
foregoing shows that there exists an isomorphism of functors
$\varphi:1_{C_1}\to GF$ such that $F$, $G$, $\varphi$ and the identity
$1_{F(C_1)}\to FG$ define an \emph{equivalence} of the category $C_1$
with the full subcategory $F(C_1)$ of $C_2$ (T, 1.2).
\end{env}

\begin{env}[8.1.6]
\label{0.8.1.6}
Let us apply Proposition (\hyperref[0.8.1.4]{8.1.4}) to the case where the
functor $F$ is $h_{X'}$, $X'$ being an arbitrary object of $C$; the map
$\beta:\operatorname{Hom}(X,X')\to\operatorname{Hom}(h_X,h_{X'})$ is here
none other than the map $w\mapsto h_w$ defined in
(\hyperref[0.8.1.2]{8.1.2}); this map being \emph{bijective},
we see, with the terminology of (\hyperref[0.8.1.5]{8.1.5}), that:
\end{env}

\begin{proposition}[8.1.7]
\label{0.8.1.7}
The canonical functor
$h:C\to\operatorname{Hom}(C^0,\mathrm{Ens})$ is fully faithful.
\end{proposition}

\begin{env}[8.1.8]
\label{0.8.1.8}
Let $F$ be a contravariant functor from $C$ to $\mathrm{Ens}$; we say that
$F$ is \emph{representable} if there exists an object $X\in C$ such that $F$
is \emph{isomorphic} to $h_X$; it follows from
(\hyperref[0.8.1.7]{8.1.7}) that the data of an $X\in C$ and of an
isomorphism of functors $g:h_X\to F$ determine $X$ up to a unique
isomorphism. Proposition (\hyperref[0.8.1.7]{8.1.7}) also means that $h$
defines an \emph{equivalence} of $C$ with the full subcategory of
$\operatorname{Hom}(C^0,\mathrm{Ens})$ consisting of the \emph{representable
contravariant functors}. It follows moreover from
(\hyperref[0.8.1.4]{8.1.4}) that giving a functorial morphism
$g:h_X\to F$ is equivalent to giving an element $\xi\in F(X)$; saying that
$g$ is an \emph{isomorphism} is equivalent, for $\xi$, to the following
condition: \emph{for every object $Y$ of $C$, the map
$v\mapsto(F(v))(\xi)$ from $\operatorname{Hom}(Y,X)$ to $F(Y)$ is
bijective}. When $\xi$ satisfies this condition, we will say that the pair
$(X,\xi)$ \emph{represents} the representable functor $F$. By abuse of
language, we will also say that the object $X\in C$ represents $F$ if there
exists $\xi\in F(X)$ such that $(X,\xi)$ represents $F$, in other words if
$h_X$ is isomorphic to $F$.

Let $F$, $F'$ be two representable contravariant functors from $C$ to
$\mathrm{Ens}$, and $h_X\to F$ and $h_{X'}\to F'$ two isomorphisms of
functors. Then it follows from (\hyperref[0.8.1.6]{8.1.6}) that there is a
canonical bijective correspondence between $\operatorname{Hom}(X,X')$ and
the set $\operatorname{Hom}(F,F')$ of functorial morphisms $F\to F'$.
\end{env}

\begin{env}[8.1.9]
\label{0.8.1.9}
\textit{Examples. I: projective limits.} The notion of representable
contravariant functor covers in particular the notion ``dual'' to
the usual notion of ``solution of a
\oldpage[0\textsubscript{III}]{8}
universal problem''. More generally, we will see that the notion of
\emph{projective limit} is a particular case of that of representable
functor. Recall that in a category $C$, a \emph{projective system} is
defined by giving a preordered set $I$, a
family $(A_\alpha)_{\alpha\in I}$ of objects of $C$, and, for every pair
of indices $(\alpha,\beta)$ such that $\alpha\leq\beta$, a morphism
$u_{\alpha\beta}:A_\beta\to A_\alpha$. A \emph{projective limit} of this
system in $C$ consists of an object $B$ of $C$ (denoted
$\varprojlim A_\alpha$) and, for each $\alpha\in I$, a morphism
$u_\alpha:B\to A_\alpha$, such that: 1\textsuperscript{o}
$u_\alpha=u_{\alpha\beta}u_\beta$ for $\alpha\leq\beta$;
2\textsuperscript{o} For every object $X$ of $C$ and every family
$(v_\alpha)_{\alpha\in I}$ of morphisms $v_\alpha:X\to A_\alpha$ such
that $v_\alpha=u_{\alpha\beta}v_\beta$ for $\alpha\leq\beta$, there exists a
unique morphism $v:X\to B$ (denoted $\varprojlim v_\alpha$) such that
$v_\alpha=u_\alpha v$ for every $\alpha\in I$ (T, 1.8). This is interpreted
as follows: the $u_{\alpha\beta}$ canonically define
maps
\[
  \overline{u}_{\alpha\beta}:\operatorname{Hom}(X,A_\beta)
  \to\operatorname{Hom}(X,A_\alpha)
\]
which define a \emph{projective system} of sets
$(\operatorname{Hom}(X,A_\alpha),\overline{u}_{\alpha\beta})$, and
$(v_\alpha)$ is by definition an element of the set
$\varprojlim_\alpha\operatorname{Hom}(X,A_\alpha)$; it is clear that
$X\mapsto\varprojlim_\alpha\operatorname{Hom}(X,A_\alpha)$ is a
\emph{contravariant functor} from $C$ to $\mathrm{Ens}$, and the existence of
the projective limit $B$ amounts to saying that
$(v_\alpha)\mapsto\varprojlim v_\alpha$ is an \emph{isomorphism} of
functors in $X$
\[
  \varprojlim_\alpha\operatorname{Hom}(X,A_\alpha)
  \xrightarrow{\sim}\operatorname{Hom}(X,B)
  \tag{8.1.9.1}\label{0.8.1.9.1}
\]
in other words, that the functor
$X\mapsto\varprojlim_\alpha\operatorname{Hom}(X,A_\alpha)$ is
\emph{representable}.
\end{env}

\begin{env}[8.1.10]
\label{0.8.1.10}
\textit{Examples. II: Final object.} Let $C$ be a category and $\{a\}$ a
set consisting of a single element. Consider the contravariant functor
$F:C\to\mathrm{Ens}$ which assigns to every object $X$ of $C$
the set $\{a\}$, and to every morphism $X\to X'$ in $C$ the unique
map $\{a\}\to\{a\}$. To say that this functor is \emph{representable}
means that there exists an object $e\in C$ such that for every $Y\in C$,
$\operatorname{Hom}(Y,e)=h_e(Y)$ \emph{consists of a single element}; we say
that $e$ is a \emph{final object} of $C$, and it is clear that two final
objects of $C$ are isomorphic (which makes it possible to define, in general
with the help of the axiom of choice, a final object of $C$, which is then
denoted $e_C$). For example, in the category $\mathrm{Ens}$, the final objects
are the sets consisting of a single element; in the category of
\emph{augmented algebras} over a field $K$ (where the morphisms are the
algebra homomorphisms compatible with the augmentations), $K$ is a final
object; in the category of \emph{$S$-preschemes} (I, 2.5.1), $S$ is a final
object.
\end{env}

\begin{env}[8.1.11]
\label{0.8.1.11}
For two objects $X$, $Y$ of a category $C$, set
$h'_X(Y)=\operatorname{Hom}(X,Y)$, and for every morphism $u:Y\to Y'$, let
$h'_X(u)$ be the map $v\mapsto uv$ from $\operatorname{Hom}(X,Y)$ to
$\operatorname{Hom}(X,Y')$; $h'_X$ is then a \emph{covariant functor}
$C\to\mathrm{Ens}$, from which we deduce, as in
(\hyperref[0.8.1.2]{8.1.2}), the definition of a canonical covariant
functor $h':C^0\to\operatorname{Hom}(C,\mathrm{Ens})$; a \emph{covariant
functor} $F$ from $C$ to $\mathrm{Ens}$, in other words an object of
$\operatorname{Hom}(C,\mathrm{Ens})$, is then said to be \emph{representable}
if there exists an object $X\in C$ (necessarily unique up to unique
isomorphism) such that $F$ is \emph{isomorphic to $h'_X$}; we leave it to the
reader to develop the considerations ``dual'' to the preceding ones for
this notion, which this time covers that of \emph{inductive limit}, and in
particular the usual notion of ``solution of a universal problem''.
\end{env}

\subsection{Algebraic structures in categories}
\label{subsection:0.8.2}

\begin{env}[8.2.1]
\label{0.8.2.1}
\oldpage[0\textsubscript{III}]{9}
Given two contravariant functors $F$, $F'$ from $C$ to
$\mathrm{Ens}$, recall that for every object $Y\in C$ we set
$(F\times F')(Y)=F(Y)\times F'(Y)$, and for every morphism $u:Y\to Y'$ in
$C$ we set $(F\times F')(u)=F(u)\times F'(u)$, which is the map
$(t,t')\mapsto((F(u))(t),(F'(u))(t'))$ from $F(Y')\times F'(Y')$ to
$F(Y)\times F'(Y)$; $F\times F':C\to\mathrm{Ens}$ is therefore a
\emph{contravariant functor} (which, moreover, is none other than the
\emph{product} of the objects $F$, $F'$ in the category
$\operatorname{Hom}(C^0,\mathrm{Ens})$). Given an object $X\in C$, we
will call an \emph{internal composition law on $X$} a functorial morphism
\[
  \gamma_X:h_X\times h_X\to h_X.
  \tag{8.2.1.1}\label{0.8.2.1.1}
\]

In other words (T, 1.2), for every object $Y\in C$, $\gamma_X(Y)$ is a
map $h_X(Y)\times h_X(Y)\to h_X(Y)$ (hence by definition an
\emph{internal composition law} on the set $h_X(Y)$) subject to the
condition that, for every morphism $u:Y\to Y'$ in $C$, the diagram
\[
  \xymatrix{
    h_X(Y')\times h_X(Y')
      \ar[rr]^{h_X(u)\times h_X(u)}\ar[dd]_{\gamma_X(Y')} &&
    h_X(Y)\times h_X(Y)\ar[dd]^{\gamma_X(Y)}\\\\
    h_X(Y')\ar[rr]_{h_X(u)} && h_X(Y)
  }
\]
is commutative; this means that for the composition laws
$\gamma_X(Y)$ and $\gamma_X(Y')$, $h_X(u)$ is a \emph{homomorphism} from
$h_X(Y')$ to $h_X(Y)$.

In the same way, given two objects $Z$, $X$ of $C$, we call an
\emph{external composition law on $X$, having $Z$ as domain of
operators}, a functorial morphism
\[
  \omega_{X,Z}:h_Z\times h_X\to h_X.
  \tag{8.2.1.2}\label{0.8.2.1.2}
\]

We see as above that for every $Y\in C$, $\omega_{X,Z}(Y)$ is an external
composition law on $h_X(Y)$, having $h_Z(Y)$ as domain of
operators, and such that for every morphism $u:Y\to Y'$, $h_X(u)$ and
$h_Z(u)$ form a \emph{di-homomorphism} from
$(h_Z(Y'),h_X(Y'))$ to $(h_Z(Y),h_X(Y))$.
\end{env}

\begin{env}[8.2.2]
\label{0.8.2.2}
Let $X'$ be a second object of $C$, and suppose given on $X'$ an internal
composition law $\gamma_{X'}$; we will say that a morphism
$w:X\to X'$ in $C$ is a \emph{homomorphism} for these composition
laws if for every $Y\in C$, $h_w(Y):h_X(Y)\to h_{X'}(Y)$ is a
homomorphism for the composition laws $\gamma_X(Y)$ and
$\gamma_{X'}(Y)$. If $X''$ is a third
\oldpage[0\textsubscript{III}]{10}
object of $C$ equipped with an internal composition law $\gamma_{X''}$, and
$w':X'\to X''$ a morphism in $C$ which is a homomorphism for
$\gamma_{X'}$ and $\gamma_{X''}$, it is clear that the morphism
$w'w:X\to X''$ is a homomorphism for the composition laws $\gamma_X$
and $\gamma_{X''}$. An isomorphism $w:X\xrightarrow{\sim}X'$ in $C$ is
called an \emph{isomorphism for the composition laws $\gamma_X$ and
$\gamma_{X'}$} if $w$ is a homomorphism for these composition laws and
if its inverse morphism $w^{-1}$ is a homomorphism for the composition
laws $\gamma_{X'}$ and $\gamma_X$.

The \emph{di-homomorphisms} for pairs of objects of $C$ equipped with
external composition laws are defined in the same way.
\end{env}

\begin{env}[8.2.3]
\label{0.8.2.3}
When an internal composition law $\gamma_X$ on an object $X\in C$ is
such that $\gamma_X(Y)$ is a \emph{group law} on $h_X(Y)$ for
\emph{every} $Y\in C$, we say that $X$, equipped with this law, is a
\emph{$C$-group} or a \emph{group object in $C$}. The
\emph{$C$-rings}, \emph{$C$-modules}, etc.\ are defined in the same way.
\end{env}

\begin{env}[8.2.4]
\label{0.8.2.4}
Suppose that the \emph{product} $X\times X$ of an object $X\in C$ with itself
exists in $C$; by definition, we then have
$h_{X\times X}=h_X\times h_X$ up to canonical isomorphism, since this is
a particular case of projective limit
(\hyperref[0.8.1.9]{8.1.9}); an internal composition law on $X$ can
therefore be regarded as a functorial morphism
$\gamma_X:h_{X\times X}\to h_X$, and therefore canonically determines
(\hyperref[0.8.1.6]{8.1.6}) an element
$c_X\in\operatorname{Hom}(X\times X,X)$ such that $h_{c_X}=\gamma_X$; in
this case, giving an internal composition law on $X$ is therefore
equivalent to giving a morphism $X\times X\to X$; when $C$ is the
category $\mathrm{Ens}$, we recover the classical notion of internal
composition law on a set. There is an analogous result for an external
composition law when the product $Z\times X$ exists in $C$.
\end{env}

\begin{env}[8.2.5]
\label{0.8.2.5}
With the preceding notation, suppose moreover that $X\times X\times X$
exists in $C$; the characterization of the product as an object representing a
functor (\hyperref[0.8.1.9]{8.1.9}) implies the existence of canonical
isomorphisms
\[
  (X\times X)\times X\xrightarrow{\sim}X\times X\times X
  \xrightarrow{\sim}X\times(X\times X)~;
\]
if we canonically identify $X\times X\times X$ with $(X\times X)\times X$,
the map $\gamma_X(Y)\times 1_{h_X(Y)}$ is identified with
$h_{c_X\times 1_X}(Y)$ for every $Y\in C$. Consequently it is equivalent to
say that for every $Y\in C$ the internal law $\gamma_X(Y)$ is
\emph{associative}, or that the diagram of maps
\[
  \xymatrix{
    h_X(Y)\times h_X(Y)\times h_X(Y)
      \ar[rr]^{\gamma_X(Y)\times 1}\ar[dd]_{1\times\gamma_X(Y)} &&
    h_X(Y)\times h_X(Y)\ar[dd]^{\gamma_X(Y)}\\\\
    h_X(Y)\times h_X(Y)\ar[rr]_{\gamma_X(Y)} && h_X(Y)
  }
\]
\oldpage[0\textsubscript{III}]{11}
is commutative, or that the diagram of morphisms
\[
  \xymatrix{
    X\times X\times X
      \ar[rr]^{c_X\times 1_X}\ar[dd]_{1_X\times c_X} &&
    X\times X\ar[dd]^{c_X}\\\\
    X\times X\ar[rr]_{c_X} && X
  }
\]
is commutative.
\end{env}

\begin{env}[8.2.6]
\label{0.8.2.6}
Under the hypotheses of (\hyperref[0.8.2.5]{8.2.5}), if we want to express
that for every $Y\in C$ the internal law $\gamma_X(Y)$ is a \emph{group
law}, we must on the one hand express that it is associative, and on the
other that there exists a map $\alpha_X(Y):h_X(Y)\to h_X(Y)$ having the
properties of the \emph{inverse} in a group; since, for every morphism
$u:Y\to Y'$ in $C$, we have seen that $h_X(u)$ must be a group
homomorphism $h_X(Y')\to h_X(Y)$, we see first that $\alpha_X:h_X\to h_X$
must be a \emph{functorial morphism}. On the other hand, we can express the
characteristic properties of the inverse $s\mapsto s^{-1}$ in a group
$G$ without bringing in the neutral element: it suffices to write that the
two composite maps
\[
  (s,t)\mapsto(s,s^{-1},t)\mapsto(s,s^{-1}t)\mapsto s(s^{-1}t),
\]
\[
  (s,t)\mapsto(s,s^{-1},t)\mapsto(s,ts^{-1})\mapsto(ts^{-1})s
\]
are equal to the second projection $(s,t)\mapsto t$ from $G\times G$ to
$G$. By virtue of (\hyperref[0.8.1.3]{8.1.3}), we have
$\alpha_X=h_{a_X}$, where $a_X\in\operatorname{Hom}(X,X)$; the first
of the preceding conditions then expresses that the composite morphism
\[
  X\times X\xrightarrow{(1_X,a_X)\times 1_X}X\times X\times X
  \xrightarrow{1_X\times c_X}X\times X\xrightarrow{c_X}X
\]
is the second projection $X\times X\to X$ in $C$, and the second condition
is expressed in the same way.
\end{env}

\begin{env}[8.2.7]
\label{0.8.2.7}
Suppose now that there exists in $C$ a \emph{final object} $e$
(\hyperref[0.8.1.10]{8.1.10}). Suppose still that $\gamma_X(Y)$ is
a group law on $h_X(Y)$ for every $Y\in C$, and denote by
$\eta_X(Y)$ the neutral element of $\gamma_X(Y)$. Since, for every morphism
$u:Y\to Y'$ in $C$, $h_X(u)$ is a group homomorphism, we have
$\eta_X(Y)=(h_X(u))(\eta_X(Y'))$; taking in particular $Y'=e$, in which case
$u$ is the unique element $\varepsilon$ of $\operatorname{Hom}(Y,e)$, we see
that the element $\eta_X(e)$ completely determines $\eta_X(Y)$ for every
$Y\in C$. Set $e_X=\eta_X(X)$, the neutral element of the group
$h_X(X)=\operatorname{Hom}(X,X)$; the commutativity of the diagram
\[
  \xymatrix{
    h_X(e)\ar[r]^{h_X(\varepsilon)}\ar[d]_{h_{e_X}(e)} &
    h_X(Y)\ar[d]^{h_{e_X}(Y)}\\
    h_X(e)\ar[r]_{h_X(\varepsilon)} & h_X(Y)
  }
\]
(cf. \hyperref[0.8.1.2]{8.1.2}) shows that, in the set $h_X(Y)$,
the map $h_{e_X}(Y)$ is none other than
\oldpage[0\textsubscript{III}]{12}
$s\mapsto\eta_X(Y)$, which transforms every element into the neutral element.
One then verifies that the fact that $\eta_X(Y)$ is the neutral element of
$\gamma_X(Y)$ for every $Y\in C$ is equivalent to saying that the composite
morphism
\[
  X\xrightarrow{(1_X,1_X)}X\times X
  \xrightarrow{1_X\times e_X}X\times X\xrightarrow{c_X}X,
\]
and the analogous one in which $1_X$ and $e_X$ are interchanged, are both
\emph{equal} to $1_X$.
\end{env}

\begin{env}[8.2.8]
\label{0.8.2.8}
One could of course easily multiply the examples of algebraic
structures in categories. The example of groups has been treated in some
detail, but in what follows we will generally leave it to the reader to
develop analogous considerations in the examples of algebraic structures
that we will encounter.
\end{env}

